\subsection{Which horizontal strips carry the doubled families?}\label{sup:lienard-product}
Proposition~\othernumber{prop:strip-product}{6.3} of the main article counts sign brackets on \(k=\lfloor(\Delta+1)/2\rfloor\) separate strips. After \(p\) vertical folds, the force \(V_p'\) has \(2^p\) odd powers and the perturbation \(F_p\) has at most \(2^pb-1\) nonconstant powers. At \(m=\lceil3\Delta/2\rceil\) terms, replicating \((V_p,F_p)\) directly gives leading coefficient \(1/4\). On the diagonal \((N,N-1)\), a terminal even fold doubles \(c_p\) to \(2c_p\) without new outer terms, yielding \(2^{p+1}\) odd force terms and at most \(2^pb-1\) even perturbation terms of degree at most \(2^{p+1}b-2\), which also raises the diagonal leading coefficient to \(1/4\). Han and Romanovski's property \(Z(n,m,k)\) concerns \(k\) odd-multiplicity cycles for sufficiently small perturbations of \(\dot y=X-\varepsilon F(y),\ \dot X=-g(y)\), with \(\deg F\le n+1,\ \deg g\le m\) and \(g(+\infty)=+\infty\) \cite{HanRomanovski}. In its first doubling branch \(n'=2n+1,\ m'=2m+1,\ k'=2k+n\). The sign-bracket realization uses \(n=s_i-1\) in this count.

The coordinate \(X=AT_k(x)\) has one regular inverse near each of the \(k\) simple zeros. There \(U=X^2/2\) and the confining vertical potential gives the outer ovals used in the proof. The intervening maxima have local form \(U=A^2/2-X^2/2\) in a different chart. A large outer bracket for \(X^2/2+V_p(y)\) therefore has no copy there by this change of variables. In the stationary construction of Proposition~\othernumber{prop:stationary-budget}{6.1} the vertical potential is odd, and reflection in \(y\) also supplies the maximum strips.

The horizontal replication in this proof works for any \(k\ge1\) whose strip degree and term count fit. The column \(\mathcal T\) in main Table~\othernumber{tab:fold}{2} fixes \(k=\lfloor N/2\rfloor\). For example, \((p,b,k)=(6,7,449)\) at \(N=1024\) and \((8,15,3841)\) at \(N=8192\) give \(1{,}150{,}338\) and \(116{,}036{,}610\) by the same formula \(2kc_p\). In both examples the term ceiling before filling is exactly \(N\), with degrees at most \(897\) and \(7681\), respectively. Any unused positions within the allowances may again be filled with small rational coefficients.
